\subsection{Motivation and Background}

Neural networks have demonstrated remarkable success in computer vision and natural language processing domains, achieving state-of-the-art performance on tasks ranging from medical diagnosis to autonomous systems. However, their application to physics-based problems reveals fundamental limitations in domain generalization. While trained models excel at interpolating within their training distributions, they struggle to maintain accuracy when encountering variations in physical environments, boundary conditions, or parameter regimes that characterize real-world physics applications.

The emerging field of Simulation Intelligence represents a convergence of scientific computing, numerical simulation, and neural network architectures~\cite{Paszke2019PyTorch,Ott2020A,Prashanthi2025Deep}. Recent advances have introduced several promising frameworks for integrating neural networks with physics solvers. Physics-Informed Neural Networks (PINNs) embed governing differential equations directly into the loss function, enabling the network to learn solutions that satisfy physical constraints~\cite{Mishra2020Estimate,Ryck2024Numerica}. Theoretical work has begun to establish generalization error bounds for PINNs in specific contexts, including ecological diffusion models and time-fractional diffusion equations~\cite{Reyes2025A,Trinh2025Converge}. Neural Ordinary Differential Equations (NODEs) parameterize continuous-time dynamics, offering a framework for learning temporal evolution operators. Universal Differential Equations (UDEs) combine mechanistic models with neural network components to capture unknown physics. These architectures have demonstrated substantial benefits including orders-of-magnitude computational speedups, reduced implementation costs, and the ability to adapt partial differential equations to conform with observational data~\cite{Ren2022Towards,Wong2025Evolutio}.

Operator learning methods extend these capabilities by learning mappings between function spaces rather than finite-dimensional vectors. Fourier Neural Operators (FNOs) leverage spectral representations to achieve resolution-invariant predictions for parametric partial differential equations~\cite{Li2020Fourier,Li2022Fourier,Wen2021UFNO}. Recent extensions address complex geometries, discontinuous coefficients, and boundary conditions through learned deformations, Walsh-Hadamard bases, and boundary-embedded architectures~\cite{Raoni'c2023Convolut,Cavallazzi2025WalshHad,Wang2024BENO}. DeepONet architectures learn operators by separately encoding input functions and query locations, demonstrating effectiveness in elastoplastic structures, magnetotelluric modeling, and financial time series~\cite{He2023Novel,Liao2024Fast,Ahmad2024DeepONet}. Physics-constrained variants incorporate conservation laws and temporal integration schemes to improve accuracy~\cite{Jnini2025Physicsc,Mandl2025PhysicsI}. Specialized operator learning frameworks address challenges including dimensionality reduction, spectral accuracy, arbitrary domain geometries, and stability for hyperbolic and Hamiltonian systems~\cite{Mandl2024Separabl,Sojitra2025FEDONet,Mousavi2025RIGNO,Xiao2024Neural,Tanaka2024Neural}. Theoretical analyses demonstrate that operator learning can mitigate the curse of dimensionality for certain PDE classes and enable data-free learning through spectral methods~\cite{Chen2023Deep,Choi2023Spectral,Rahman2024Pretrain,Zhang2025PhysicsI}.

Despite these architectural innovations, a critical gap remains: current neural network solutions for physics problems lack evidence of robust test performance across the entire range of physically realizable inputs. The literature contains numerous demonstrations comparing neural network solutions to numerical simulators on selected examples within narrow parameter regimes. However, numerical solutions to differential equations are expected to work universally across all environments, realistic physical relationships, and boundary conditions. Trained deep neural network models consistently fail to generalize when encountering even modest variations in environmental conditions, material properties, or forcing terms. This limitation is particularly acute in domains where the input space is high-dimensional and the physically plausible region forms a complex manifold within that space.

Generalization theory for deep learning has advanced significantly in computer vision and natural language processing contexts. Research has explored complexity measures, sharpness-aware minimization, data augmentation strategies, and architectural constraints to improve out-of-distribution performance~\cite{Neyshabur2017Explorin,Natekar2020Represen,Vakanski2021EVALUATI,Zhang2025Position}. Sharpness-Aware Minimization (SAM) and its variants seek flat minima in the loss landscape, demonstrating improved generalization across diverse tasks including federated learning, time series forecasting, and scale-invariant learning~\cite{Foret2020Sharpnes,Kwon2021ASAM,Qu2022Generali,Andriushchenko2022Towards,Ilbert2024SAMforme,Li2024Friendly,Luo2025Explicit}. Data augmentation techniques tailored for domain generalization have shown promise in reinforcement learning, medical image segmentation, and contrastive learning frameworks~\cite{Raileanu2020Automati,Pinneri2023Equivari,Su2022Rethinki,Jiang2025CbDA,Kumar2025RandSali}. Architectural innovations including random CNN structures and representation-based complexity measures provide additional tools for predicting and improving generalization~\cite{\u015awiderski2022Random,Rajagopal2020Predicti,Litman2026A}. However, these techniques have not been systematically evaluated in physics domains where the structure of the problem space differs fundamentally from image or text distributions.

Physics applications present unique characteristics that distinguish them from traditional deep learning domains. In underwater acoustics, transmission loss prediction depends on complex interactions between sound speed profiles, bathymetry, and propagation physics governed by parabolic equation models~\cite{Lalita2024Underwat,Tu2023A,Wang2024Ocean,Agustinus2024ANALYSIS,Kayser2023Validity}. Machine learning approaches have been applied to acoustic field modeling and sound speed profile prediction, but generalization across diverse oceanographic conditions remains challenging~\cite{Sun2025Endtoend,Mccarthy2023ReducedO,Huang2021Rapid,Ahmed2021Machine}. Radio frequency propagation in the atmosphere exhibits similar complexity, with refractivity gradients, ducting conditions, and atmospheric stratification creating highly variable propagation environments~\cite{Bidigare2025DualPola,McQuaid2025Overview,Pastore2025Identify,Bidigare2023Data}. Modified refractivity gradient models and numerical weather prediction systems provide ground truth for training, but neural network surrogates must generalize across geographic regions, seasonal variations, and synoptic weather patterns~\cite{Kalu2018Determin,Saleem2016Statisti,Parada2023A,Javeed2018A,Pu2022Modified,Emmanuel2020Geostati,Rabinovich2025Using,Nguyen2023Machine}. Structural health monitoring using guided waves requires detecting and localizing damage across varying environmental conditions, material properties, and sensor configurations~\cite{Zheng2023A,Cao2019Probabil,Yang2024Modeling,Liu2022An,Chen2021Fatigue,Paialunga2023Damage}. Time-frequency analysis techniques and unsupervised learning methods address temperature compensation and outlier detection, but achieving robust performance across the full operational envelope remains an open challenge~\cite{Yang2022Guidelin,Yang2021Longterm,Tabjula2021Outlier,Yue2020A,Zhou2025A,Singh2025Hybrid,Han2025Timefreq,Zhang2025TimeFreq,Tansel2022NEW,Liu2020A,Huang2021A}.

Recent advances in contrastive learning and vision foundation models offer potential pathways for improving physics domain generalization. Contrastive pretraining methods learn representations by maximizing agreement between augmented views of the same example while minimizing agreement between different examples~\cite{Yang2022Chinese,Lin2023PMCCLIP,Brannon2023ConGraT,Zhou2025Contrast}. Applications span medical imaging, computational pathology, and multimodal vision-language tasks~\cite{Mih\u0103i\u0163\u01032025CLIPXRad,Ahmed2025CLIPRL,Ilyas2025CNNLSTM,Jiao2025ImageTex}. Various contrastive loss formulations including angular margin, signed quadratic, and inverse contrastive losses have been developed for specific domains~\cite{Memon2023Robustne,Dehghan2024CCLDTI,Li2023Revisiti,Raichur2024Bayesian,Guo2023A,Hayashi2021Deep,Akash2021Learning,Choi2020AMCLoss}. Vision foundation models pretrained on large-scale datasets demonstrate strong transfer learning capabilities across medical imaging, semantic segmentation, and domain generalization tasks~\cite{Ryu2025Visionla,Pai2025Vision,Li2025MedDINOv,Huang2025LoftUp,Zhang2025Mamba,Zhao2025FisherTu,Chen2023Intern,Mai2024On}. However, the applicability of these representation learning techniques to physics domains, where the input-output relationships are governed by differential equations rather than statistical correlations, remains unexplored.

The fundamental gap addressed by this research is the absence of systematic studies on generalization in physics domains. While theoretical generalization bounds exist for specific PINN formulations~\cite{Mishra2020Estimate,Reyes2025A,Trinh2025Converge}, no prior work has investigated which properties and methods from mature deep learning domains transfer to physics applications, nor has research developed physics-specific generalization techniques that exploit domain structure, symmetries, and conservation laws. Numerical solutions to differential equations are expected to work universally across all physically realizable conditions, but current neural network approaches fail to meet this standard. This project addresses this gap by conducting foundational research to enable neural networks to generalize across broad spectra of physical environments and boundary conditions in Navy-relevant physics tasks.